\documentclass[11pt,a4paper]{article}
\usepackage[T1]{fontenc}
\usepackage{XCharter,amsmath,amssymb}
\usepackage[margin=25mm]{geometry}
\usepackage{hyperref}
\title{A waiting-policy clarification for the recursive auction example}
\author{Pathfinder experiment record}
\date{28 September 2026}
\begin{document}
\maketitle
This corrigendum concerns the September 25 repeat--reinjection pilot,
reinjection thread Q3P5, ledger entry 19 and its supporting matched-state
construction. The archived output is unchanged. This clarification is not
an independently reviewed paper or an input to the next experiment.

\paragraph{The gap.}
After an initial bid of 2, the buyer of value 3 might cross the standing
ask of 2.50 on a later invitation. If it does so before the seller of
cost 1 arrives, the event that this seller is invited within six draws
no longer guarantees a trade with surplus 2. Merely requiring agents to
avoid losses does not establish the written lower bound.

\paragraph{An explicit sufficient continuation.}
Retain the stated eight active traders and their values, the two hidden
cost assignments, and the six independent uniform invitations among
eligible traders. Fix other traders' behaviour: they hold their quotes,
except that seller $S_2$ crosses a strictly profitable standing bid when
invited. Traders who have transacted exit the round.
The acting buyer is allowed to choose its continuation.
In state A, an informed comparator posts 2 and leaves it standing for all
six draws. Until $S_2$ is invited there is no trade; all eight traders
remain eligible. Consequently its expected added welfare is exactly
$2r$, where $r=1-(7/8)^6$.

Immediate crossing yields 0.50 in either state and at most another 0.25
in state A, even if later profitable trades are permitted. In state B,
after a noncrossing first action no seller can profitably cross a bid
below 2.50. A profitable trade requires the buyer to be invited again;
its surplus is at most 0.50. Thus that branch has expected welfare at
most $0.50r$, while immediate crossing yields 0.50.

Under an equal prior, let $p$ be the probability that an observation-only
policy crosses immediately. Its state-informed regret is bounded below by
\[
\frac12\bigl[p(2r-0.75)+(1-p)0.50(1-r)\bigr]
\geq \frac12\min\{2r-0.75,\;0.50(1-r)\}
=0.1121988296508789\ldots.
\]
The first term compares with a feasible informed waiting policy, not an
assumed optimal policy; the second allows the uninformed buyer to adapt
after later observations. This is sufficient for the regret inequality.
It does not mean that every noncrossing continuation earns $2r$.

\paragraph{Scope.}
This repairs the existence construction by specifying behaviour, not by
demonstrating that the stipulated behaviour is an equilibrium or typical
of the source experiment. The hidden-state prior is illustrative. No
frequency estimate, simulator-equivalence test, or causal attribution of
observed market losses follows. The original conditional-on-a-final-actor
repeat example remains a related, distinct construction.
\end{document}
